\section{Concluding Remarks}
\label{sec:conclusion}

We have proven that for \kl{hereditary classes} of graphs of \kl{bounded
clique-width}, being \kl{$\forall$-well-quasi-ordered} is a very robust notion,
that enjoys many structural characterisations, and where most conjectures from
the literature on labelled well-quasi-orderings hold true. We consider these
results as a first step towards understanding the landscape of labelled
well-quasi-orderings of graphs and other relational structures.
Note that we proved a very strong dichotomy result for hereditary
classes of graphs of bounded clique-width: either they are not
\kl{$2$-well-quasi-ordered} or the \kl{induced subgraph} relation on them is
simpler than the \kl{gap-embedding} relation on trees. We conjecture that this
dichotomy extends to all hereditary classes of graphs: there is nothing beyond
(nested) trees.
We strongly believe that our techniques can be adapted to the case
of finite relational structures over a finite relational signature. Another
interesting direction would be to consider non-hereditary classes of graphs of
bounded clique-width. In this case, we also believe that our techniques can be
adapted.

A fundamental limitation of our work is that we assume classes to be of
\kl{bounded clique-width}. We believe that classes of graphs that are
\kl{$2$-well-quasi-ordered} are necessarily of \kl{bounded clique-width} (see
\cref{cwqo:conj}), and a first step would be to show that every
such class is \emph{monadically dependent}, a notion originating from 
model theory \cite{baldwin85}
that has been recently investigated in finite model theory 
with the belief that it characterises tractable model checking for 
first-order formulas (see for instance \cite{dreier24} and \cite{dreier2024flipbreakability}).
We actually conjecture a stronger statement, that does not involve labelling
the graphs of the class.

\begin{conjecture}
  \label{conj:wqo-mon-dep}
  Every \kl{hereditary class} of graphs that is 
  \kl{well-quasi-ordered} by the \kl{induced subgraph} relation is \emph{monadically dependent}.
\end{conjecture}

Even in the case of classes having \kl{bounded clique-width}, there are 
several unanswered conjectures that are refinement of our results.
For instance, can we
\kl{existentially interpret} all paths in every \kl{hereditary class} of graphs
that is not \kl{$\forall$-well-quasi-ordered} (and of \kl{bounded
clique-width}). An argument similar to \cref{lem:transduction-not-wqo} shows
that if a class \kl{existentially interprets} all paths, then it is not
\kl{$2$-well-quasi-ordered}, but the converse direction remains open.

Similarly, we focused on classes that are \kl{$2$-well-quasi-ordered}
as our simplest possible labelling. However, there are other 
ways to add information to the graphs: one can distinguish $n$ vertices
(adding $n$ constants), or one can assume that the labelling 
itself is ordered by $a \leq b$ ; adapting the notion of embedding in each case. 
It is clear that a class that is 
\kl{$2$-well-quasi-ordered} will be \kl{well-quasi-ordered} when
finitely many constants are added to distinguish nodes 
\cite[Lemma 5.2]{braunfeld21}, and that it will still be 
\kl{well-quasi-ordered} when \kl{freely labelling} the class
with two comparable labels $a \leq b$. We believe that the converse holds.

\begin{conjecture}
  Every \kl{hereditary class} of graphs that is 
  \kl{well-quasi-ordered} by the \kl{induced subgraph} relation 
  with \emph{two} additional constants
  is \kl{$2$-well-quasi-ordered}. 
\end{conjecture}